\section{Proposed Methodology}
\label{sec:method}

The framework is organised in three phases, together with a set of cross-cutting
properties that apply to all of them. Phase~1 extracts multi-modal signals from the raw
data, Phase~2 fuses those signals conditionally on the identified market regime and
converts the fused view into portfolio weights, and Phase~3 evaluates the result and
attributes it to its components. The level of detail below is deliberate, so that the
study can be reproduced independently from this description alone.

\subsection{Phase 1: Multi-Modal Signal Extraction}

\subsubsection{Input Layer}

The architectural framework initiates with the processing of market data for the target
securities, structured as open-high-low-close-volume (OHLCV) series:

\begin{itemize}
\item \textbf{Market data.} Daily open, high, low, close and volume records for the
instruments in the study universe, adjusted for splits and dividends, obtained from
Yahoo Finance through the \texttt{yfinance} interface. The download window is
1~January 2007 to 27~August 2026. The end date is not a modelling choice: it is every
session available at the date the analysis was run, so the sample boundary carries no
information about the results.
\item \textbf{News and social media.} Not used. No news article, social media post,
corporate filing, transcript or analyst note is read at any point in this study. The
reason, and the measure adopted in place of a textual one, are given in
Section~\ref{sec:sentiment}.
\item \textbf{Macroeconomic data.} Not used as an input. Published macroeconomic and
sentiment series are employed once, in Section~\ref{sec:sentval}, and there only as
external validators that never enter the framework.
\end{itemize}

These raw inputs undergo a preprocessing pipeline involving cleansing, temporal
normalisation and synchronisation to construct a cohesive, time-aligned dataset for
subsequent analytical stages.

The primary study universe consists of ten exchange-traded funds selected to span the
major liquid asset classes available to a diversified allocator: US large capitalisation
equity (SPY), US large capitalisation growth equity (QQQ), US small capitalisation
equity (IWM), developed international equity (EFA), emerging market equity (EEM),
long-maturity US Treasuries (TLT), intermediate-maturity US Treasuries (IEF),
investment-grade corporate credit (LQD), gold (GLD) and US real estate (VNQ). Breadth
matters for the research question. A framework whose stated purpose is to redirect risk
across market states can only be tested where there is somewhere else for risk to go; a
single-sector universe offers no such destination, a point returned to in
Section~\ref{sec:megacap}.

Indicator warm-up and the 252-session covariance estimation window consume the first
year, so the evaluated out-of-sample period runs from January 2008 to August 2026 and
covers the global financial crisis, the 2011 sovereign debt episode, the 2020 pandemic
crash, the 2022 inflation-driven bear market and the subsequent recovery. Sessions from
1~January 2024 onward are treated as a holdout. That block postdates the design of the
framework and every specification decision reported here; none of it was examined while
the model was being built. It enters the full-sample results for completeness and is
reported separately in Section~\ref{sec:holdout}. A secondary universe of ten
mega-capitalisation US equities over 2018--2026 is carried through the same pipeline and
reported in Section~\ref{sec:megacap}.

\begin{figure*}[t]
\centering
\includegraphics[width=0.70\textwidth]{images/fig_architecture.png}
\caption{Proposed Methodology. Phase~1 constructs causal, cross-sectionally standardised
features; Phase~2 identifies the market state, fuses the two composites and separates
relative composition from total risk; Phase~3 evaluates the result and ablates its
components.}
\label{fig:architecture}
\end{figure*}

\subsubsection{Technical Indicators Engine}

This component transforms raw market data into quantitative trading signals through three
primary analytical categories. Nine features are computed for each asset on each session.
All are causal: every value at session $t$ uses only information available at the close
of session $t$.

\begin{itemize}
\item \textbf{Price-derived metrics.} The relative strength index (RSI) identifies
potential overbought or oversold asset conditions. The moving average convergence
divergence (MACD) line captures trend momentum shifts. The stochastic oscillator locates
the close within its recent trading range.
\item \textbf{Volume-analysis tools.} On-balance volume (OBV) assesses momentum through
cumulative volume flow, entered as a twenty-session change normalised by average volume,
while the money flow index (MFI) combines price and volume into a single participation
measure. Relative volume measures turnover against its own recent average.
\item \textbf{Volatility estimators.} Realised volatility measures the statistical
dispersion of asset returns over a specified look-back period, and volume instability
measures the dispersion of daily volume changes.
\end{itemize}

No Bollinger band, average true range or volume profile construction is used; the nine
features of Table~\ref{tab:features} are the complete input set.

Working mechanism: a sliding-window methodology facilitates all calculations, with
parameters calibrated for the horizons given in Table~\ref{tab:features}. No other input
enters the framework.

\begin{table}[!ht]
\centering
\caption{Input features. $C_t$, $H_t$, $L_t$ and $V_t$ denote close, high, low and
volume. All features are computed per asset, then cross-sectionally standardised.}
\label{tab:features}
\footnotesize
\begin{tabular}{p{2.0cm}p{1.8cm}p{5.2cm}p{1.5cm}}
\hline
\textbf{Feature} & \textbf{Composite} & \textbf{Definition} & \textbf{Window} \\
\hline
RSI & Technical & Relative strength index, centred: $\mathrm{RSI}_t - 50$ & 14 \\
\hline
MACD & Technical & Moving average convergence divergence line, scaled by price: $\mathrm{MACD}_t / C_t$ & 12, 26, 9 \\
\hline
Stochastic & Technical & Stochastic oscillator \%K, centred: $\%K_t - 50$ & 14, smooth 3 \\
\hline
MFI & Technical & Money flow index, centred: $\mathrm{MFI}_t - 50$ & 14 \\
\hline
OBV flow & Technical & Twenty-session change in on-balance volume, normalised by average volume: $(\mathrm{OBV}_t - \mathrm{OBV}_{t-20}) / \overline{V}_{t,20}$ & 20 \\
\hline
Return & Sentiment & Five-session simple return: $C_t / C_{t-5} - 1$ & 5 \\
\hline
Relative volume & Sentiment & $V_t / \overline{V}_{t,20}$ & 20 \\
\hline
Realised volatility & Sentiment & Annualised standard deviation of daily returns, entered with a negative sign & 20 \\
\hline
Volume instability & Sentiment & Standard deviation of daily volume changes, entered with a negative sign & 20 \\
\hline
\end{tabular}
\end{table}

The final output constitutes a normalised signal matrix in which each column corresponds
to a technical feature scaled to a uniform numerical range. Each feature is standardised
cross-sectionally on every session, so that a value expresses how an asset ranks against
its peers that day rather than against its own history. Writing $x_{i,t}$ for the raw
feature of asset $i$ and $N$ for the number of assets,
\begin{equation}
z_{i,t} = \frac{x_{i,t} - \bar{x}_t}{s_t}, \qquad
\bar{x}_t = \frac{1}{N}\sum_{j} x_{j,t}, \qquad
s_t = \mathrm{sd}_j(x_{j,t}).
\end{equation}
Cross-sectional standardisation is used rather than time-series standardisation because
the allocation problem is a relative one: the optimiser must decide which assets to
prefer today, not whether today is unusual by historical standards. It also removes the
common market factor from the scores, so that a market-wide move does not appear as a
signal on every asset simultaneously.

The standardised features are then averaged into two composites. The technical composite
is the equally weighted mean of the five technical features. The sentiment composite is
the equally weighted mean of the four sentiment features, taken with the signs given in
Table~\ref{tab:features}, so that recent strength and active participation raise the
score while realised volatility and unstable volume lower it.

\subsubsection{Sentiment Analysis Pipeline}
\label{sec:sentiment}

This component transforms market data into a quantitative sentiment indicator. The
sentiment measure used here is market-derived. It is the equally weighted mean of the
four features marked \emph{Sentiment} in Table~\ref{tab:features}: a five-session return,
turnover measured against its own recent average, realised volatility and the instability
of volume, the last two entered with a negative sign. The composite therefore rises when
participation broadens and prices are advancing calmly, and falls when trading becomes
disorderly. Every input is a price or a volume, and every one is causal.

Because readers of work in this area reasonably ask which text model was used, it is
worth being categorical: none. Neither a domain-specific transformer model nor a
lexicon-based method in the manner of \citet{loughran11} is applied, for the
straightforward reason that no text is processed for either to act upon. There is no
named entity recognition stage, because there are no entities to recognise.

That is a design decision forced by data availability rather than a preference. Textual
measurement needs a time-stamped corpus covering the traded instruments, and for broad
exchange-traded funds no such corpus exists at instrument level---a fund holding several
hundred securities has no news flow of its own. The alternative adopted here is the
market-derived tradition described in Section~\ref{sec:lit}, in which turnover and
dispersion stand in for the survey or the newspaper. Turnover is among the canonical
components of the investor-sentiment index of \citet{bakerwurgler07}.

Adopting a proxy imposes an obligation to demonstrate that the proxy behaves like the
thing it stands in for, and that obligation is discharged rather than assumed.
Section~\ref{sec:sentval} correlates the composite against four independently published
series, none of which enters the framework at any stage, and reports how the identified
market states separate on it. What that exercise cannot establish is that the composite
predicts returns; Section~\ref{sec:ic} measures that separately, and reports a result
that is unfavourable to the framework. The two questions are distinct and are kept
distinct throughout.

\subsubsection{Regime Detection Module}

This component identifies latent, non-observable market states through a statistical
clustering methodology.

\begin{itemize}
\item \textbf{Probabilistic state modelling.} The system employs a Gaussian mixture model
(GMM) to segment market data into distinct probabilistic clusters, following the
tradition of state-dependent modelling established by \citet{hamilton89}. The clustering
uses two features, each computed as a cross-sectional average across the universe: the
twenty-session rolling standard deviation of daily returns, and the ten-session rolling
mean of daily returns. The first captures the volatility of the environment, the second
its direction. Both are standardised before fitting.
\item \textbf{Volatility clustering.} A specific function of the module is to isolate and
characterise periods in which either high or low volatility persists, effectively
partitioning the market timeline into volatility-based regimes.
\item \textbf{Qualitative labelling.} The identified states are labelled by ascending
fitted mean volatility, which yields a calm state, a transitional state and a stress
state, and Section~\ref{sec:regimefit} reports their realised characteristics.
\end{itemize}

The procedure follows a sequential pipeline: initial feature extraction from market data,
application of the probabilistic clustering algorithm, and the qualitative labelling of
the identified states. The final outputs comprise a time series of regime assignments and
an estimated state transition matrix.

The mixture uses $K = 3$ components with full covariance matrices, three random restarts
and a fixed seed. It is refitted at every rebalance date on an expanding window
containing only sessions strictly earlier than that date, with a minimum of 252
observations; the state of the current session is then obtained by prediction from that
fit. No observation ever contributes to the model that labels it.

One implementation detail is essential and is easy to overlook. The component indices
returned by a mixture fit are arbitrary and permute freely between refits, so labels from
consecutive fits are not comparable unless they are anchored. Components are therefore
sorted in ascending order of fitted mean volatility before use, which makes regime~0 the
calm state, regime~1 the transitional state and regime~2 the stress state consistently
across every refit. Omitting this step does not produce an error; it produces a silently
scrambled regime series and, if the resulting labels are compared across fits, a spurious
appearance of instability.

The choice $K = 3$ is not selected by information criterion. As
Section~\ref{sec:regimefit} reports, both the Akaike information criterion (AIC) and the
Bayesian information criterion (BIC) continue to fall as components are added over the
range tested, which is the usual behaviour of mixture models on persistent financial
series and reflects their tendency to absorb heavy tails with additional components. The
choice is instead made on persistence and interpretability: three components yield states
that last long enough to act on and that admit a clear economic reading, whereas finer
partitions fragment into short-lived states that a monthly rebalancing schedule cannot
exploit. This is a modelling judgement and is reported as such.

\subsection{Phase 2: Decision-Theoretic Fusion Framework}

\subsubsection{Signal Confidence Weighting}

This mechanism executes a synthesis of predictive signals contingent upon the prevailing
market regime. Its operational logic is detailed below.

\begin{itemize}
\item \textbf{Regime-dependent weights.} Signal influence is modulated by the identified
market state. Technical indicators are accorded greater authority during stable,
low-volatility regimes. Conversely, during periods characterised by elevated uncertainty,
sentiment-based proxies assume a dominant role in the composite forecast. The weights
$(w^T, w^M)$ are
$(0.7, 0.3)$ in the calm state, $(0.5, 0.5)$ in the transitional state and $(0.3, 0.7)$
in the stress state.
\item \textbf{Uncertainty quantification.} The framework quantifies the reliability of
the combined view per asset by the agreement between its two composites. The confidence
term down-weights assets whose composites disagree:
\begin{equation}
c_{i,t} = \mathrm{clip}\!\left(0.5 + 0.5\,\rho^{(63)}_{i,t}, \; 0.1, \; 1.0\right),
\end{equation}
where $\rho^{(63)}_{i,t}$ is the trailing 63-session correlation between the two
composites for asset $i$. Agreement between independent views is treated as evidence;
disagreement is treated as noise and shrinks the score toward zero.
\item \textbf{Fixed schedule, tested rather than learned.} The weight schedule is
specified in advance from the hypothesis stated in Section~\ref{sec:intro}---that
trend-following information is more reliable in quiet trending markets, while attention
and stress information becomes relatively more informative when markets are
disturbed---and is then held fixed. It is not re-estimated from realised performance, and
no Bayesian updating is applied. Fitting the schedule to the sample it is evaluated on
would make the resulting improvement uninterpretable; instead,
Section~\ref{sec:ic} measures the conditional efficacy the schedule presumes, and reports
whether that presumption holds.
\end{itemize}

\subsubsection{Regime-Dependent Signal Fusion}

This mechanism synthesises the weighted predictive inputs into a unified forecast, with
its operational logic conditioned upon the prevailing market regime.

\begin{itemize}
\item \textbf{Conditional combination.} The methodology applies a distinct combination
rule to each identified market state, on the hypothesis that the predictive efficacy and
optimal weighting of signals are regime-contingent.
\item \textbf{Correlation structure.} The fusion process explicitly maintains the
dependency structure across different assets, ensuring that the combined forecast reflects
existing inter-asset correlations. The fused scores are passed to an optimiser carrying
the full estimated covariance matrix, so the combined view is expressed as a portfolio
rather than as a set of independent per-asset bets.
\item \textbf{Uncertainty propagation.} The resulting composite forecasts incorporate the
estimation uncertainty present within the constituent input signals, providing a more
robust output. This operates through the confidence term of the preceding subsection,
which shrinks the score of any asset whose two composites disagree.
\end{itemize}

Operational framework: the core engine executes a multiplication of the regime-specific
signal weight vector against the matrix of individual composites, adjusted by the
per-asset confidence term. Let $T_{i,t}$ and $M_{i,t}$ denote the standardised technical
and sentiment composites for asset $i$, and let $r_t \in \{0,1,2\}$ be the regime. The
fused score is
\begin{equation}
z_{i,t} = c_{i,t} \left[ w^{T}_{r_t} T_{i,t} + w^{M}_{r_t} M_{i,t} \right].
\end{equation}

The fused score is then converted into an expected excess return using the standard
information-ratio scaling of \citet{grinold00}, in which an assumed information
coefficient (IC) places the score on a return scale rather than leaving it as a
dimensionless quantity:
\begin{equation}
\alpha_{i,t} = \mathrm{IC} \cdot \sigma_{i,t} \cdot z_{i,t},
\label{eq:alpha}
\end{equation}
where $\sigma_{i,t}$ is the annualised volatility of asset $i$ from the current
covariance estimate and $\mathrm{IC} = 0.05$.

Equation~(\ref{eq:alpha}) is not a cosmetic step. If a raw indicator is passed to a
mean-variance objective directly, the units of the linear term and of the quadratic risk
term are unrelated, and their ratio---not the information in the signal---decides the
solution. An indicator such as on-balance volume takes values of order $10^{8}$ to
$10^{9}$, whereas annualised variances are of order $10^{-2}$; the linear term then
dominates by ten orders of magnitude and the optimiser degenerates to a corner solution,
holding a single asset regardless of the risk model. Section~\ref{sec:alloc} documents
this failure mode explicitly, because it is the source of a behaviour visible in the
originally submitted results.

\subsubsection{Portfolio Optimization Engine}

This component translates fused signals into executable allocations. Allocation proceeds
in two steps, which separate two decisions that are often conflated: which assets to hold
relative to one another, and how much total risk to carry.

\begin{itemize}
\item \textbf{Shrinkage covariance.} $\Sigma_t$ is estimated by Ledoit--Wolf shrinkage
\citep{ledoit04} on the trailing 252 daily returns and annualised by a factor of 252.
Shrinkage is used because a sample covariance matrix estimated from 252 observations on
ten assets is poorly conditioned, and a mean-variance optimiser will concentrate weight
precisely on the directions in which the estimate is least reliable.
\item \textbf{Risk-budgeting allocation.} Total exposure is set by the identified regime
rather than by the optimiser, so that the composition decision and the risk-level
decision are made separately and can be evaluated separately.
\item \textbf{Transaction cost optimisation.} The objective carries an explicit linear
penalty on turnover, so that the cost of trading enters the decision rather than being
deducted afterwards.
\end{itemize}

\paragraph{Step 1: relative allocation.} The risky sleeve solves a quadratic program with
a linear transaction-cost penalty:
\begin{align}
\max_{w} \quad & \alpha_t^{\top} w - \lambda \, w^{\top} \Sigma_t w
- \kappa \left\| w - w_{t-1} \right\|_1 \\
\text{s.t.} \quad & w \geq 0, \qquad w \leq w_{\max}, \qquad \mathbf{1}^{\top} w = 1,
\end{align}
where $\Sigma_t$ is the annualised covariance matrix, $\lambda = 2$ is the risk aversion
coefficient, $w_{\max} = 0.25$ caps any single position at one quarter of the sleeve, and
$w_{t-1}$ is the weight vector carried in from the previous rebalance. The turnover
penalty $\kappa = c \cdot 252 / h$ annualises a one-way cost of $c = 10$ basis points
over a rebalancing interval of $h = 21$ sessions, so that the cost of trading is expressed
on the same annual basis as the return and risk terms. The problem is convex and is
solved with CVXPY using the Clarabel interior-point solver.

The equality constraint $\mathbf{1}^{\top} w = 1$ is what makes this step purely relative.
It fixes the sleeve to be fully invested, so the optimiser can only express a view about
composition, never about how much risk to take. That second decision is made explicitly
in step~2 instead of emerging as a by-product of the constraint set.

\paragraph{Step 2: regime risk budget.} The sleeve is then scaled by a factor that depends
only on the regime:
\begin{equation}
w^{\ast}_t = k_{r_t} \, \tilde{w}_t, \qquad
k = \{\,\text{calm}: 1.00, \;\; \text{transitional}: 0.75, \;\; \text{stress}: 0.50\,\},
\label{eq:budget}
\end{equation}
with the residual $1 - k_{r_t}$ held in cash, which is assumed to earn zero. Scaling is
one-sided: the framework de-risks in disturbed states but never levers up in calm ones,
so any performance difference is attributable to risk reduction rather than to leverage.

The budget in Equation~(\ref{eq:budget}) is expressed as a fraction of the sleeve's own
risk rather than as an absolute volatility target, and this choice is consequential. An
absolute target de-risks only when the sleeve's estimated volatility exceeds the target
level, so its effect depends on the volatility of the universe it is applied to. Targets
calibrated for a concentrated equity portfolio---of the order of 10 to 20 per cent
annualised---are simply never reached by a diversified ten-asset sleeve whose own realised
volatility is close to 7~per~cent. Under those targets the constraint binds in one
rebalance out of 224 in this study, and the regime layer is inert: it is present in the
specification, consumes computation, and does nothing. A relative budget is scale-free by
construction and remains active in any universe. Section~\ref{sec:ablation} reports both
variants and the difference between them.

\subsection{Phase 3: Performance Evaluation \& Feedback}

\subsubsection{Backtesting Framework}

The proposed methodology undergoes historical validation through a structured simulation
protocol designed to emulate real-world trading conditions.

\begin{itemize}
\item \textbf{Walk-forward validation.} The implementation utilises a rolling-origin
technique, sequentially advancing a 252-session model estimation window followed by a
21-session out-of-sample evaluation period. Between rebalances positions are held and
allowed to drift with realised returns, so that reported turnover reflects trading
actually undertaken rather than a notional daily reweighting. Every estimated quantity
uses only data preceding the decision it informs:
the covariance matrix from the trailing 252 sessions before the rebalance date, the
mixture model from an expanding window strictly earlier than that date, and all features
from causal rolling windows.
\item \textbf{Out-of-sample testing.} A strict temporal partition is enforced between
parameter calibration and performance assessment phases, to prevent statistical
overfitting and ensure robust results. No parameter is tuned on the evaluation period; the
configuration reported as primary is the one specified above, and
Section~\ref{sec:robust} reports the full grid of alternatives rather than the best cell
of it. The 665 sessions from January 2024 onward postdate every specification decision
and are reported separately in Section~\ref{sec:holdout}.
\item \textbf{Regime-aware analysis.} Strategy efficacy is disaggregated and evaluated
contingent upon the market state identified by the regime detection module, allowing
diagnosis of performance drivers by state rather than in aggregate only.
\end{itemize}

The framework executes a historical simulation that incorporates realistic market
frictions. One-way transaction costs of 10 basis points are charged on the traded
fraction $\|w^{\ast}_t - w_{t-1}\|_1$ on the rebalance session;
Section~\ref{sec:cost} varies this assumption from 0 to 25 basis points.

\subsubsection{Performance Attribution}

This analytical procedure deconstructs aggregate returns to isolate the fundamental
sources of value generation, employing a three-part approach.

\begin{itemize}
\item \textbf{Regime-based decomposition.} The methodology disaggregates performance
contributions across distinct market environments, isolating the returns attributable to
the calm, transitional and stress states identified by the regime detection module.
\item \textbf{Signal contribution analysis.} The framework quantifies the marginal
performance impact originating from each signal category, thereby identifying which
inputs were the primary drivers of the result. Component contributions are isolated by
ablation rather than by assertion: the regime layer is removed while holding average
exposure fixed, removed while holding full exposure, and the signal layer is removed
entirely. Holding average exposure fixed is what distinguishes the effect of timing risk
from the effect of simply holding less, and without it the two are not separable. The
information coefficient of each composite against next-session cross-sectional returns is
measured separately, overall and by regime.
\item \textbf{Risk-adjusted metrics.} Evaluation incorporates risk-adjusted metrics
including the Sharpe and Calmar ratios, alongside compound annual growth rate (CAGR),
annualised volatility and maximum drawdown (Max DD); the Sharpe ratio is computed at a
zero risk-free rate. These are supplemented with bootstrap-based statistical significance
tests to distinguish robust performance from random outcomes. Statistical inference on
daily excess returns uses a paired $t$-test, a Newey--West $t$-statistic with ten lags to
accommodate autocorrelation and heteroskedasticity \citep{newey87}, and a stationary block
bootstrap with 21-session blocks and 2{,}000 replications \citep{politis94} for interval
estimates of the Sharpe ratio and of the Sharpe difference.
\end{itemize}

Three benchmarks are used. The first is an equal-weighted portfolio of the same ten
assets, rebalanced daily and charged no transaction costs, which is deliberately a
generous comparator \citep{demiguel09}. The second is an inverse-volatility portfolio
using trailing 60-session standard deviations, also costless, which isolates whether the
framework adds anything beyond elementary risk weighting. The third rescales the
equal-weighted portfolio to the strategy's realised volatility, which removes the
mechanical advantage a lower-risk portfolio enjoys in ratio comparisons.

\subsubsection{Model Refinement Loop}

The framework incorporates a recursive re-estimation mechanism that keeps its statistical
components current as the sample advances. What that mechanism does, and what it
deliberately does not do, are both stated here.

\begin{itemize}
\item \textbf{Regime model updates.} The underlying statistical model for state
identification is systematically recalibrated to reflect structural shifts in market
dynamics, preventing reliance on obsolete regime definitions. The mixture is refitted at
each of the 224 rebalance dates on an expanding window of all sessions strictly earlier
than that date.
\item \textbf{Risk model updates.} The shrinkage covariance matrix is re-estimated at
each rebalance from the trailing 252 sessions, and the shrinkage intensity is re-derived
from that window rather than held constant.
\item \textbf{Parameters that are \emph{not} re-optimised.} The regime weight schedule,
the risk budget ladder, the assumed information coefficient, the risk aversion
coefficient and the position cap are fixed ex ante and never adapted to realised
performance. No automated parameter search, signal decommissioning rule or
human-in-the-loop override operates inside the evaluation period. This is a restriction
rather than a feature, and it is adopted because a loop that re-optimises on realised
performance inside the evaluation window makes the resulting backtest uninterpretable
\citep{harvey15}.
\end{itemize}

\subsection{Cross-Cutting Features}

\subsubsection{Temporal Consistency}

The framework maintains chronological integrity across all analytical processes through
three mechanisms.

\begin{itemize}
\item All incoming data streams are aligned to a unified temporal index, eliminating
artefacts arising from asynchronous sampling frequencies.
\item Causal relationships are explicitly parameterised through the lag structure of
every rolling window: each feature is computed from data strictly preceding the session
it describes, and every model is fitted on data strictly preceding the decision it
informs. No observation contributes to the model that labels it, and no quantity
estimated at any date uses information from a later one.
\item Regime labels are anchored across refits by sorting mixture components on fitted
mean volatility, so that a state identified at one rebalance means the same thing at the
next. Without this, the component indices permute freely between fits and the regime
series becomes incomparable over time.
\end{itemize}

\subsubsection{Risk Management Integration}

The framework incorporates a layered risk protocol designed to operate within real-world
constraints.

\begin{itemize}
\item The system enforces predetermined boundaries on individual position sizes: weights
are non-negative and capped at $w_{\max} = 0.25$ of the risky sleeve. These constraints
prevent excessive concentration and ensure adherence to diversification principles, and
they matter because an unconstrained mean-variance solution concentrates precisely where
its estimated inputs are least reliable.
\item De-risking operates at the aggregate portfolio level through the regime budget of
Equation~(\ref{eq:budget}), with the residual held in cash. No discretionary stop-loss
rule is used, and none is needed: the budget is a standing rule applied at every
rebalance rather than a trigger fired after a loss has been realised.
\item Turnover is penalised inside the objective rather than charged after the fact, so
that a position is taken only when the expected improvement exceeds the cost of trading
into it. Section~\ref{sec:cost} reports the response to costs of 0 to 25 basis points.
\end{itemize}

\subsubsection{Explainability and Transparency}

The framework incorporates specific mechanisms to ensure analytical clarity throughout
the decision process.

\begin{itemize}
\item Every allocation decision can be directly attributed to its originating predictive
signal and the prevailing market regime, providing complete traceability. The weight
decomposes into a relative weight produced by the optimiser from the fused score and a
scalar exposure set by the regime; because the two are applied in sequence rather than
jointly, the contribution of each is separately observable, and
Section~\ref{sec:ablation} measures both.
\item The principal performance claims are presented with their associated confidence
intervals, quantifying the degree of statistical certainty for each output: Sharpe ratios
and Sharpe differences carry block-bootstrap intervals.
\item The architecture maintains a comprehensive, timestamped record of all inputs, model
states and decision logic, creating a verifiable audit trail suitable for systematic
strategy validation. The analysis writes its results to machine-readable artefacts, and
every numerical claim in Section~\ref{sec:results} is checked programmatically against
those artefacts rather than transcribed by hand.
\end{itemize}

This architecture, shown in Figure~\ref{fig:architecture}, creates a system in which
portfolio decisions follow from the identified market state and the fused signal by an
explicit and inspectable route, while maintaining causal statistical foundations and
practical implementability.
